% Book 4, Year 2, chapter 28 -- Search and games.
% Provenance (never printed): sources/book4/BRIEFS.md, chapter 28.
\chapter{Search and Games}\label{ch:b2:search-and-games}

The sliding-tile puzzle has eight numbered tiles in a $3 \times 3$ frame and one empty cell; a move slides a tile into
the empty cell, and the goal is to put the tiles in order. A program can try every sequence of moves, but how many it
has to try depends on how it chooses:

\omcode{code/bachelor-2/28-search-and-games/listings/eight_puzzle.py}{38}{47}{Five scrambled puzzles solved three ways,
counting the states each search expands}\label{lst:b2:search-and-games:hook}

\omoutput{code/bachelor-2/28-search-and-games/outputs/eight_puzzle.txt}{The same shortest solutions, found after very
different amounts of work.}\label{out:b2:search-and-games:hook}

All three searches return a shortest solution, and the third expands fifty times fewer states than the first on the
hardest puzzle. The difference is an estimate of the distance left, used carefully enough that the answer stays
optimal. The second half of the chapter adds an opponent: in noughts and crosses, or in a race between a robot and a
rival through the town, each player chooses in turn, and a program can still compute who wins, and how.

\begin{recall}
The Year 1 volume: breadth-first and depth-first search; Dijkstra's algorithm for shortest paths with non-negative
weights, and its proof; priority queues kept as binary heaps; backtracking. The grade 10 volume: graphs, paths and
trees.
\end{recall}

\section{State spaces}\label{sec:b2:search-and-games:state-spaces}

\begin{definition}[State space, search problem, goal test]\label{def:b2:search-and-games:state-space}
A \emph{search problem}\index{search problem} is given by a set of \emph{states}, a start state, a \emph{goal
test}\index{goal test} that says whether a state is a goal, and for each state its \emph{successors}, each reached by an
action of cost $c > 0$. The states and the actions form a directed graph, the \emph{state space}\index{state space}; a
solution is a path from the start to a goal, and an optimal solution one of least total cost.
\end{definition}

The \omterm{def:b2:search-and-games:state-space}{state space} is usually never built: the 8-puzzle has $9!/2 = 181{,}440$ reachable states, chess far more than any
memory holds. The program knows only the start and a function that lists successors; the graph is \emph{implicit}, and
a search builds the part it visits.

\begin{omfigure}
\begin{tikzpicture}[x=1cm, y=1cm,
  pics/puzzle/.style={code={
    \foreach \v [count=\i from 0] in {#1} {
      \pgfmathtruncatemacro{\r}{\i/3}\pgfmathtruncatemacro{\c}{mod(\i,3)}
      \ifnum\v=0 \node[draw, fill=black!12, minimum size=4mm, inner sep=0pt] at (\c*0.4, -\r*0.4) {};
      \else \node[draw, minimum size=4mm, inner sep=0pt, font=\scriptsize] at (\c*0.4, -\r*0.4) {\v}; \fi}}}]
\pic at (4.6, 0) {puzzle={1,2,3,4,0,6,7,5,8}};
\node[font=\small, right] at (5.6, -0.4) {start, $h = 2$};
\pic at (0, -2.4) {puzzle={1,0,3,4,2,6,7,5,8}};
\pic at (3, -2.4) {puzzle={1,2,3,0,4,6,7,5,8}};
\pic at (6.2, -2.4) {puzzle={1,2,3,4,6,0,7,5,8}};
\pic at (9.2, -2.4) {puzzle={1,2,3,4,5,6,7,0,8}};
\foreach \x/\h in {0/3, 3/3, 6.2/3, 9.2/1} \node[font=\small] at (\x + 0.4, -3.6) {$h = \h$};
\pic at (9.2, -4.9) {puzzle={1,2,3,4,5,6,7,8,0}};
\node[font=\small, right] at (10.2, -5.3) {goal};
\foreach \x in {0, 3, 6.2} \draw[->, >=Stealth] (5, -0.95) -- (\x + 0.4, -2.2);
\draw[->, >=Stealth, very thick, omDef] (5, -0.95) -- (9.6, -2.2);
\draw[->, >=Stealth, very thick, omDef] (9.6, -3.85) -- (9.6, -4.7);
\end{tikzpicture}
\par\smallskip{\small Part of the 8-puzzle's \omterm{def:b2:search-and-games:state-space}{state space}. The start has four successors, one per tile that can slide;
each is labelled with its Manhattan distance $h$, the sum over the tiles of their distances to their places. The
thick path is the solution of two moves; $h$ points straight at it.}
\end{omfigure}

\omcode{code/bachelor-2/28-search-and-games/listings/state_space.py}{15}{26}{Breadth-first search in an implicit graph,
counting the states expanded}\label{lst:b2:search-and-games:bfs}

Breadth-first search expands the states in order of depth, so it finds a solution with fewest actions. If each state has
at most $b$ successors and the shallowest goal is at depth $d$, it may expand $1 + b + \dots + b^d = O(b^d)$ states, and
keeps them all in memory. Depth-first search keeps only the current path, but may go down forever, or find a long
solution first.

\begin{definition}[Iterative deepening]\label{def:b2:search-and-games:iddfs}
\emph{Iterative deepening}\index{iterative deepening} runs a depth-first search limited to depth $0$, then $1$, then $2$,
and so on, until a limit at which it finds a goal.
\end{definition}

It finds a solution with fewest actions, like breadth-first search, while keeping only a path of $d + 1$ states. Repeating
the shallow levels costs little: with $b \geq 2$ the last level holds at least half of all the states up to it, and the
total work is $O(b^d)$ again (\cref{exo:b2:search-and-games:3,exo:b2:search-and-games:4}).

\needspace{10\baselineskip}
\omcode{code/bachelor-2/28-search-and-games/listings/state_space.py}{29}{51}{Depth-limited search, which never revisits a
state of the current path, and \omterm{def:b2:search-and-games:iddfs}{iterative deepening}}\label{lst:b2:search-and-games:iddfs}

\section{Uniform-cost search}\label{sec:b2:search-and-games:ucs}

When actions have different costs, as roads have different lengths, fewest actions is not cheapest.

\begin{definition}[Uniform-cost search]\label{def:b2:search-and-games:ucs}
\emph{Uniform-cost search}\index{uniform-cost search} keeps a priority queue of states ordered by $g(s)$, the cost of the
cheapest path to $s$ found so far. It repeatedly removes a state of smallest $g$, marks it \emph{expanded}, stops if it is
a goal, and otherwise updates $g$ for its successors not yet expanded.
\end{definition}

It is Dijkstra's algorithm of the Year 1 volume, run on an implicit graph and stopped at the first goal.

\begin{theorem}[Uniform-cost search]\label{thm:b2:search-and-games:ucs}
With positive action costs, \omterm{def:b2:search-and-games:ucs}{uniform-cost search} expands states in order of non-decreasing distance from the start; when
it expands a state $s$, $g(s)$ is the distance to $s$. In particular the first goal expanded is reached by an optimal
solution.
\end{theorem}

\begin{proof}
This is the invariant of Dijkstra's algorithm, proved in the Year 1 volume: if a state $s$ were expanded with $g(s)$
larger than its distance, the first state $u$ of an optimal path to $s$ that is not yet expanded would have, from its
expanded predecessor, $g(u)$ equal to its distance, at most the distance of $s$, which is less than $g(s)$, so $u$ would
have been removed from the queue first. A goal expanded first is at the least distance among the goals: any closer goal
would have been expanded before.
\end{proof}

\section{Heuristics and A*}\label{sec:b2:search-and-games:astar}

\omterm{def:b2:search-and-games:ucs}{Uniform-cost search} spreads in every direction. Most problems come with knowledge of the direction to go: the
straight-line distance to the destination, or how far each tile is from its place.

\begin{definition}[Heuristic, admissible, consistent]\label{def:b2:search-and-games:heuristic}
A \emph{heuristic}\index{heuristic} is a function $h \geq 0$ on states that estimates the cost of reaching a goal, with
$h = 0$ on goals. Write $h^*(s)$ for the true least cost from $s$ to a goal. The heuristic is
\emph{admissible}\index{admissible heuristic} if $h(s) \leq h^*(s)$ for every $s$: it never overestimates. It is
\emph{consistent}\index{consistent heuristic} if $h(s) \leq c + h(t)$ for every action from $s$ to $t$ of cost $c$.
\end{definition}

\begin{proposition}[Consistent implies admissible]\label{prop:b2:search-and-games:consistent-admissible}
A \omterm{def:b2:search-and-games:heuristic}{consistent heuristic} is \omterm{def:b2:search-and-games:heuristic}{admissible}.
\end{proposition}

\begin{proof}
\Cref{exo:b2:search-and-games:8}: follow an optimal path from $s$ to a goal and add the inequalities.
\end{proof}

\begin{method}[A*]\label{met:b2:search-and-games:astar}\index{A* search}
\omterm{def:b2:search-and-games:ucs}{Uniform-cost search} in which the priority of a state $s$ is $f(s) = g(s) + h(s)$, the cost so far plus the estimate of the
cost left.
\end{method}

\omcode{code/bachelor-2/28-search-and-games/listings/state_space.py}{54}{73}{A*: \omterm{def:b2:search-and-games:ucs}{uniform-cost search} when $h = 0$}\label{lst:b2:search-and-games:astar}

\begin{theorem}[Optimality of A*]\label{thm:b2:search-and-games:astar-optimal}
If $h$ is \omterm{def:b2:search-and-games:heuristic}{consistent}, A* as in \cref{lst:b2:search-and-games:astar}, which never expands a state twice, returns an optimal
solution, and it expands the states in order of non-decreasing $f$.
\end{theorem}

\begin{proof}
Give each action from $s$ to $t$ the new cost $c' = c - h(s) + h(t)$, which is $\geq 0$ by consistency. Along any path from
the start $s_0$ to a state $t$, the new costs add up to the old cost minus $h(s_0)$ plus $h(t)$: the middle terms cancel.
So the new distance to $t$ is $g(t) - h(s_0) + h(t)$, the priority of A* minus a constant. A* is therefore \omterm{def:b2:search-and-games:ucs}{uniform-cost
search} on the new costs, and by \cref{thm:b2:search-and-games:ucs} (whose proof needs only non-negative costs) it
expands states in order of new distance, that is of $f$, with $g$ exact. On goals $h = 0$, so the order of new distances
among goals is that of true distances, and the first goal expanded is an optimal one.
\end{proof}

If $h$ is only \omterm{def:b2:search-and-games:heuristic}{admissible}, A* that may expand a state again, when it finds a cheaper path to it, is still optimal; the
listing, which never expands a state twice, may then return a longer path. A \omterm{def:b2:search-and-games:heuristic}{heuristic} that overestimates loses the
guarantee altogether (\cref{exo:b2:search-and-games:2}).

For the 8-puzzle, the number of misplaced tiles and the Manhattan distance are both \omterm{def:b2:search-and-games:heuristic}{consistent}: one move changes one
tile's position by one cell, so it changes either \omterm{def:b2:search-and-games:heuristic}{heuristic} by at most $1$ (\cref{exo:b2:search-and-games:9}). The
Manhattan distance is never smaller than the number of misplaced tiles, and \cref{out:b2:search-and-games:hook} shows
what that buys: the closer $h$ is to $h^*$, the fewer states A* expands.

\begin{omfigure}
\begin{tikzpicture}
\begin{semilogyaxis}[width=0.72\linewidth, height=5.4cm, xlabel={length of the shortest solution (moves)},
  ylabel={states expanded}, xtick={8,10,16,20,24}, legend pos=outer north east, legend style={font=\small},
  ymin=8, ymax=300000]
\addplot[mark=*, black!60] table[x=depth, y=bfs] {code/bachelor-2/28-search-and-games/figures/eight_puzzle.dat};
\addplot[mark=square*, omThm] table[x=depth, y=misplaced] {code/bachelor-2/28-search-and-games/figures/eight_puzzle.dat};
\addplot[mark=triangle*, omDef] table[x=depth, y=manhattan] {code/bachelor-2/28-search-and-games/figures/eight_puzzle.dat};
\legend{breadth-first, A* misplaced, A* Manhattan}
\end{semilogyaxis}
\end{tikzpicture}
\par\smallskip{\small The five puzzles of \cref{out:b2:search-and-games:hook} on a logarithmic scale. All three counts
grow exponentially with the depth; a better \omterm{def:b2:search-and-games:heuristic}{heuristic} lowers the rate of growth, not only the level.}
\end{omfigure}

On the town's map, the straight-line distance between two places is \omterm{def:b2:search-and-games:heuristic}{consistent}, by the triangle inequality, and no road
is shorter than the straight line between its ends.

\omoutput{code/bachelor-2/28-search-and-games/outputs/town_astar.txt}{Three routes on the town map: the same lengths, and
a third to a half of the expansions.}\label{out:b2:search-and-games:town}

\begin{omfigure}
\begin{tikzpicture}[x=1.45cm, y=1.45cm,
  base/.style={circle, draw, minimum size=4.6mm, inner sep=0pt, font=\tiny},
  other/.style={base, draw=black!40, text=black!60},
  dijkstra/.style={base, fill=black!15},
  astar/.style={base, fill=omDef!35, thick},
  road/.style={black!30},
  route/.style={omDef, very thick}]
% written by town_fig.py from the town map; do not edit
\node[dijkstra] (p1) at (4.315, 1.430) {1};
\node[astar] (p2) at (2.301, 2.855) {2};
\node[dijkstra] (p3) at (4.990, 0.673) {3};
\node[dijkstra] (p4) at (2.764, 0.604) {4};
\node[dijkstra] (p5) at (0.033, 1.480) {5};
\node[astar] (p6) at (1.313, 2.962) {6};
\node[dijkstra] (p7) at (2.871, 1.398) {7};
\node[dijkstra] (p8) at (1.345, 4.875) {8};
\node[astar] (p9) at (2.982, 3.224) {9};
\node[astar] (p10) at (2.876, 2.818) {10};
\node[astar] (p11) at (2.005, 3.350) {11};
\node[dijkstra] (p12) at (0.636, 3.840) {12};
\node[astar] (p13) at (0.645, 2.687) {13};
\node[astar] (p14) at (1.217, 0.753) {14};
\node[dijkstra] (p15) at (1.888, 1.720) {15};
\node[other] (p16) at (3.035, 0.020) {16};
\node[dijkstra] (p17) at (2.877, 4.267) {17};
\node[dijkstra] (p18) at (5.279, 3.534) {18};
\node[astar] (p19) at (1.601, 1.434) {19};
\node[astar] (p20) at (1.084, 2.154) {20};
\node[dijkstra] (p21) at (5.080, 3.036) {21};
\node[dijkstra] (p22) at (3.557, 4.776) {22};
\node[astar] (p23) at (3.536, 2.174) {23};
\node[dijkstra] (p24) at (4.252, 0.161) {24};
\node[dijkstra] (p25) at (1.795, 4.897) {25};
\node[dijkstra] (p26) at (5.657, 1.999) {26};
\node[dijkstra] (p27) at (3.833, 4.317) {27};
\node[dijkstra] (p28) at (3.870, 3.243) {28};
\node[dijkstra] (p29) at (3.951, 3.779) {29};
\node[dijkstra] (p30) at (0.953, 4.093) {30};
\draw[road] (p1) -- (p3);
\draw[road] (p1) -- (p23);
\draw[road] (p1) -- (p24);
\draw[road] (p1) -- (p26);
\draw[road] (p2) -- (p9);
\draw[road] (p2) -- (p10);
\draw[route] (p2) -- (p11);
\draw[road] (p3) -- (p24);
\draw[road] (p3) -- (p26);
\draw[road] (p4) -- (p7);
\draw[road] (p4) -- (p15);
\draw[road] (p4) -- (p16);
\draw[road] (p5) -- (p13);
\draw[road] (p5) -- (p14);
\draw[road] (p5) -- (p20);
\draw[route] (p6) -- (p11);
\draw[road] (p6) -- (p12);
\draw[road] (p6) -- (p13);
\draw[route] (p6) -- (p20);
\draw[road] (p7) -- (p15);
\draw[road] (p7) -- (p16);
\draw[road] (p7) -- (p23);
\draw[road] (p8) -- (p12);
\draw[road] (p8) -- (p25);
\draw[road] (p8) -- (p30);
\draw[road] (p9) -- (p10);
\draw[road] (p9) -- (p11);
\draw[road] (p9) -- (p17);
\draw[road] (p9) -- (p28);
\draw[road] (p10) -- (p23);
\draw[road] (p12) -- (p13);
\draw[road] (p12) -- (p30);
\draw[road] (p13) -- (p20);
\draw[road] (p14) -- (p15);
\draw[route] (p14) -- (p19);
\draw[road] (p15) -- (p19);
\draw[road] (p15) -- (p20);
\draw[road] (p16) -- (p24);
\draw[road] (p17) -- (p22);
\draw[road] (p17) -- (p25);
\draw[road] (p17) -- (p27);
\draw[road] (p18) -- (p21);
\draw[road] (p18) -- (p28);
\draw[road] (p18) -- (p29);
\draw[route] (p19) -- (p20);
\draw[road] (p21) -- (p26);
\draw[road] (p21) -- (p28);
\draw[road] (p22) -- (p27);
\draw[road] (p22) -- (p29);
\draw[road] (p25) -- (p30);
\draw[road] (p27) -- (p28);
\draw[road] (p27) -- (p29);
\draw[road] (p28) -- (p29);

\end{tikzpicture}
\par\smallskip{\small The town's places at their coordinates, from the Station (2) to the Old Mill (14). \omterm{def:b2:search-and-games:ucs}{Uniform-cost
search} expands 29 places, shaded grey or coloured, spreading in every direction; A* expands only the 10 coloured ones,
which lie towards the goal. The thick line is the route both return. Drawn by \code{town_fig.py} from the same search.}
\end{omfigure}

\begin{history}
% ledger: b2:hnr-1968
Peter Hart, Nils Nilsson and Bertram Raphael published A* in 1968, with the proof that an \omterm{def:b2:search-and-games:heuristic}{admissible} estimate keeps the
answer optimal.
\end{history}

\section{Two-player reachability games}\label{sec:b2:search-and-games:games}

\omimg{book4/ai/b2-search-and-games-board-game-cafe.jpg}{0.6\linewidth}{Two players at a café table. In a game like
this one, each move is chosen by the player whose turn it is, and the question is no longer which path exists but which
outcome a player can force.}

\begin{definition}[Reachability game, position, game graph]\label{def:b2:search-and-games:game}
A \emph{reachability game}\index{reachability game} is played by two players, $0$ and $1$, on a directed graph, the
\emph{game graph}\index{game graph}, whose vertices, the \emph{positions}\index{position (game)}, each belong to one
player, with a set $T$ of target positions. A token starts on a position; the \omterm{def:b2:rust-ownership:ownership}{owner} of the position where the token
stands moves it along an edge. Player $p$ wins if the token reaches $T$; otherwise, including when the play goes on
forever or stops at a position with no edge, player $p$ does not win.
\end{definition}

\begin{definition}[Strategy, winning strategy, winning position]\label{def:b2:search-and-games:strategy}
A \emph{strategy}\index{strategy (game)} for player $p$ chooses a successor for each position of $p$ given the play so
far; it is \emph{memoryless}\index{memoryless strategy} if the choice depends only on the current position. It is
\emph{winning}\index{winning strategy} from a position if every play from there in which $p$ follows it reaches $T$,
whatever the other player does; that position is then a \emph{winning position}\index{winning position} for $p$.
\end{definition}

\begin{definition}[Attractor]\label{def:b2:search-and-games:attractor}
The \emph{attractor}\index{attractor} of $T$ for player $p$ is built in layers: layer $0$ is $T$; a position not yet in
it joins at layer $i + 1$ if it belongs to $p$ and has a successor in layer $i$ or before, or belongs to the other player
and all its successors (at least one) are in layers $\leq i$. The attractor is the union of the layers.
\end{definition}

\omcode{code/bachelor-2/28-search-and-games/listings/games.py}{7}{30}{The \omterm{def:b2:search-and-games:attractor}{attractor} computed backwards, with a counter of
the successors left for each position of the opponent}\label{lst:b2:search-and-games:attractor}

\begin{theorem}[Attractor]\label{thm:b2:search-and-games:attractor}
The \omterm{def:b2:search-and-games:attractor}{attractor} of $T$ for player $p$ is exactly the set of \omterm{def:b2:search-and-games:game}{positions} from which $p$ has a \omterm{def:b2:search-and-games:strategy}{winning strategy}. On it, the
\omterm{def:b2:search-and-games:strategy}{memoryless strategy} that moves from each position of $p$ to a successor of a lower layer is \omterm{def:b2:search-and-games:strategy}{winning}. The listing computes
the \omterm{def:b2:search-and-games:attractor}{attractor} and this \omterm{def:b2:search-and-games:strategy}{strategy} in time $O(|V| + |E|)$.
\end{theorem}

\begin{proof}
From a position of layer $i > 0$, if it belongs to $p$, the \omterm{def:b2:search-and-games:strategy}{strategy} moves to a lower layer; if it belongs to the other
player, every move leads to a lower layer. So the layer decreases at each move and the token reaches layer $0$, which is
$T$, within $i$ moves.

Conversely, let $u$ be outside the \omterm{def:b2:search-and-games:attractor}{attractor}. Then $u \notin T$. If $u$ belongs to $p$, all its successors are outside,
since one inside would have put $u$ in. If $u$ belongs to the other player, it has a successor outside, or no successor at
all, since otherwise $u$ would have joined. The other player can therefore always keep the token outside the \omterm{def:b2:search-and-games:attractor}{attractor}:
play never reaches $T$, whatever $p$ does, and $p$ has no \omterm{def:b2:search-and-games:strategy}{winning strategy} from $u$.

The listing processes each position once from the queue and looks at each edge once backwards, decrementing a counter;
the running time is $O(|V| + |E|)$ (\cref{exo:b2:search-and-games:12}).
\end{proof}

The theorem also says that memory is useless in \omterm{def:b2:search-and-games:game}{reachability games}: if $p$ can win at all, $p$ can win by looking only at
the current position. In the game of matches, 21 matches lie on the table, each player takes 1, 2 or 3 in turn, and
whoever takes the last match wins. A position is the number left and the player to move.

\omcode{code/bachelor-2/28-search-and-games/listings/attractor.py}{4}{16}{The game of matches as a \omterm{def:b2:search-and-games:game}{game graph}; A wins on
reaching ``no match left, B to move''}\label{lst:b2:search-and-games:matches}

\omoutput{code/bachelor-2/28-search-and-games/outputs/attractor.txt}{The \omterm{def:b2:search-and-games:game}{positions} won by the player to move, the first
move from 21, and a small game in which some \omterm{def:b2:search-and-games:game}{positions} are won by neither.}\label{out:b2:search-and-games:attractor}

The player to move loses exactly with a multiple of 4: whatever they take, the other takes the complement to 4. From 21,
take 1. In the small game of the output, each player has a target: player $0$ wants $e$, player $1$ wants $h$. Neither can
force anything from $d$, $f$ and $g$: from $d$, player $1$ goes to $f$, and $f$ and $g$ send the token back and forth for
ever. These \omterm{def:b2:search-and-games:game}{positions} are \emph{draws}.

\begin{omfigure}
\begin{tikzpicture}[x=1.8cm, y=1.3cm, p0/.style={circle, draw, thick, minimum size=7mm, inner sep=0pt},
  p1/.style={rectangle, draw, thick, minimum size=6.5mm, inner sep=0pt}, e/.style={->, >=Stealth},
  w0/.style={fill=omDef!30}, w1/.style={fill=omThm!30}, lay/.style={font=\scriptsize, above=3pt}]
\node[p0, w0] (a) at (0, 1) {$a$}; \node[lay] at (a.north) {2};
\node[p1, w1] (b) at (0, -0.2) {$b$}; \node[font=\scriptsize, below=3pt] at (b.south) {1};
\node[p0, w0] (c) at (1.2, 1) {$c$}; \node[lay] at (c.north) {1};
\node[p1] (d) at (2.4, 1) {$d$};
\node[p0, w0] (e) at (2.4, 2.1) {$e$}; \node[font=\scriptsize, left=3pt] at (e.west) {0};
\node[p1] (f) at (3.6, 1) {$f$};
\node[p0] (g) at (3.6, -0.2) {$g$};
\node[p1, w1] (h) at (1.2, -0.2) {$h$}; \node[font=\scriptsize, left=3pt] at (h.south west) {0};
\draw[e] (a) to[bend left=15] (b); \draw[e] (b) to[bend left=15] (a);
\draw[e, very thick, omDef] (a) -- (c); \draw[e, very thick, omDef] (c) -- (e);
\draw[e] (c) -- (d); \draw[e] (b) -- (h);
\draw[e] (d) -- (e); \draw[e] (d) -- (f);
\draw[e] (f) to[bend left=15] (g); \draw[e] (g) to[bend left=15] (f); \draw[e] (g) -- (h);
\draw[e] (e) edge[loop right] (e); \draw[e] (h) edge[loop below] (h);
\end{tikzpicture}
\par\smallskip{\small The small game of \cref{out:b2:search-and-games:attractor}: circles belong to player $0$, squares
to player $1$. Blue: the \omterm{def:b2:search-and-games:attractor}{attractor} of $e$ for player $0$, with its layers and the \omterm{def:b2:search-and-games:strategy}{strategy} in thick arrows. Green: the
\omterm{def:b2:search-and-games:attractor}{attractor} of $h$ for player $1$. From the white \omterm{def:b2:search-and-games:game}{positions} neither player can force a win.}
\end{omfigure}

\section{Minimax and alpha--beta}\label{sec:b2:search-and-games:minimax}

The \omterm{def:b2:search-and-games:attractor}{attractor} needs the whole \omterm{def:b2:search-and-games:game}{game graph}. For a game whose \omterm{def:b2:search-and-games:game}{positions} are too many, a program looks a few moves ahead and
estimates the rest.

\begin{definition}[Game tree, minimax value, evaluation function]\label{def:b2:search-and-games:minimax}
In a game where two players alternate, \textsc{max} and \textsc{min}, the \emph{game tree}\index{game tree} of a position
has that position at its root and the \omterm{def:b2:search-and-games:game}{positions} after each move as children. The leaves carry a value for \textsc{max}:
the result of a finished game, or the value of an \emph{evaluation function}\index{evaluation function} at the depth
where the search stops. The \emph{minimax value}\index{minimax value} of a node is its leaf value, the largest value of
its children if \textsc{max} is to move, and the smallest if \textsc{min} is.
\end{definition}

The \omterm{def:b2:search-and-games:minimax}{minimax value} is what \textsc{max} can guarantee against any play of \textsc{min}, and \textsc{min} against any play
of \textsc{max}, in the tree as given; the move to play is towards a child of the root with the root's value.

\begin{definition}[Alpha--beta pruning]\label{def:b2:search-and-games:alpha-beta}
\emph{Alpha--beta pruning}\index{alpha--beta pruning} computes the \omterm{def:b2:search-and-games:minimax}{minimax value} with two bounds passed down the tree:
$\alpha$, the value \textsc{max} is already sure of on the path from the root, and $\beta$, the value \textsc{min} is
already sure of. As soon as a node's value is known to be outside $(\alpha, \beta)$, its remaining children are skipped:
the player above has a better choice elsewhere.
\end{definition}

\omcode{code/bachelor-2/28-search-and-games/listings/alphabeta.py}{15}{35}{Alpha--beta, counting the leaves it
evaluates}\label{lst:b2:search-and-games:alphabeta}

\omoutput{code/bachelor-2/28-search-and-games/outputs/minimax_demo.txt}{A tree of depth 3 with three moves at the root and
two below: its \omterm{def:b2:search-and-games:minimax}{minimax values}, and the leaves that alpha--beta evaluates.}\label{out:b2:search-and-games:demo}

\begin{omfigure}
\begin{forest}
for tree={font=\small, l sep=10pt, s sep=5pt, inner sep=1.5pt, minimum size=5.5mm},
mx/.style={draw, regular polygon, regular polygon sides=3, inner sep=0.5pt, minimum size=7mm},
mn/.style={draw, regular polygon, regular polygon sides=3, shape border rotate=180, inner sep=0.5pt, minimum size=7mm},
cut/.style={draw=black!30, text=black!35, edge={dashed, black!35}}
[5, mx
  [5, mn [5, mx [3] [5]] [{$\geq 6$}, mx [6] [9, cut]]]
  [{$\leq 2$}, mn [2, mx [1] [2]] [7, mx, cut [0, cut] [7, cut]]]
  [{$\leq 2$}, mn [2, mx [2] [1]] [8, mx, cut [8, cut] [4, cut]]]
]
\end{forest}
\par\smallskip{\small The tree of \cref{out:b2:search-and-games:demo}: triangles pointing up are \textsc{max} nodes, down
\textsc{min} nodes. After the first subtree the root is sure of $5$. In the second subtree, \textsc{min} can already force
$2 < 5$ after two leaves, so its second child is never looked at; the same in the third. Dashed: the 5 leaves skipped.
The bounds $\geq 6$ and $\leq 2$ are all alpha--beta knows of those nodes.}
\end{omfigure}

\begin{theorem}[Alpha--beta]\label{thm:b2:search-and-games:alpha-beta}
Let $v$ be the \omterm{def:b2:search-and-games:minimax}{minimax value} of a node and $r$ the result of alpha--beta on it with bounds $\alpha < \beta$. If $\alpha <
v < \beta$, then $r = v$; if $v \leq \alpha$, then $r \leq \alpha$; if $v \geq \beta$, then $r \geq \beta$. In particular,
with $\alpha = -\infty$ and $\beta = +\infty$, alpha--beta returns the \omterm{def:b2:search-and-games:minimax}{minimax value}.
\end{theorem}

\begin{proof}
\Cref{exo:b2:search-and-games:10}, by induction on the height of the node.
\end{proof}

How much is skipped depends on the order of the children: if the best move comes first, its value sets the tightest
bound at once.

\begin{proposition}[The best case]\label{prop:b2:search-and-games:best-case}
On a uniform tree of branching $b$ and depth $d$ whose children are always ordered best first, alpha--beta evaluates
exactly $b^{\lceil d/2 \rceil} + b^{\lfloor d/2 \rfloor} - 1$ of the $b^d$ leaves, when the leaf values are distinct.
\end{proposition}

\admitted

A leaf is evaluated exactly when, along its path from the root, all the choices at the even depths are first children,
or all the choices at the odd depths are; the first kind number $b^{\lfloor d/2 \rfloor}$, the second $b^{\lceil d/2
\rceil}$, and one leaf, the leftmost, is of both kinds. The characterisation is the analysis of Knuth and Moore.

\omoutput{code/bachelor-2/28-search-and-games/outputs/alphabeta.txt}{Random uniform trees: all leaves, alpha--beta in
random order, and alpha--beta with the children sorted best first.}\label{out:b2:search-and-games:alphabeta}

With perfect ordering the count is exactly the formula, about $2b^{d/2}$: in the same time, alpha--beta searches twice as
deep as minimax. In random order it still skips most leaves. Real programs approach the best case by trying first the
moves that were best in a shallower search.

\begin{history}
% ledger: b2:von-neumann-1928
% ledger: b2:knuth-moore-1975
John von Neumann published the theory of games of \omterm{def:b2:search-and-games:strategy}{strategy}, with the minimax theorem, in 1928. Donald Knuth and Ronald
Moore published their analysis of \omterm{def:b2:search-and-games:alpha-beta}{alpha--beta pruning}, with the best case above, in 1975.
\end{history}

\section{Noughts and crosses, solved}\label{sec:b2:search-and-games:tictactoe}

Noughts and crosses is small enough to search completely. A position is a string of nine cells; $x$ moves first.

\omcode{code/bachelor-2/28-search-and-games/listings/tictactoe.py}{32}{51}{The value of a position by minimax with
memoisation, and the size of the full \omterm{def:b2:search-and-games:minimax}{game tree}}\label{lst:b2:search-and-games:tictactoe}

\omoutput{code/bachelor-2/28-search-and-games/outputs/tictactoe.txt}{The game solved: a draw with best play.}\label{out:b2:search-and-games:tictactoe}

The full tree has 549,946 nodes, but only 5,478 distinct \omterm{def:b2:search-and-games:game}{positions}: the same position is reached by many orders of the
same moves, and the memo computes each once. The value $0$ says that neither player can force a win: the game is a draw
if both play well. The last line shows a mistake that loses: after a cross in the centre, a nought on an edge lets the
crosses force a win. With the eight symmetries of the square the \omterm{def:b2:search-and-games:game}{positions} drop to 765
(\cref{exo:b2:search-and-games:11}).

\section{Exercises}\label{sec:b2:search-and-games:exercises}

\begin{exercise}[$\star$]\label{exo:b2:search-and-games:1}
Run breadth-first search on the town's roads from the Bus Depot by hand, taking the neighbours of each place in the order
of the road table, and list the places by depth. Check with \code{exo_bfs_order.py}.
\end{exercise}

\begin{exercise}[$\star$]\label{exo:b2:search-and-games:2}
Use three times the straight-line distance as the \omterm{def:b2:search-and-games:heuristic}{heuristic} of A* on the town map. Find a pair of places for which the
route returned is longer than the shortest one, and explain which step of the proof of
\cref{thm:b2:search-and-games:astar-optimal} fails. Check with \code{exo_inadmissible.py}.
\end{exercise}

\begin{exercise}[$\star$]\label{exo:b2:search-and-games:3}
Compare \omterm{def:b2:search-and-games:iddfs}{iterative deepening} with breadth-first search on scrambled 8-puzzles: states expanded, and states kept in memory.
Use \code{exo_iddfs.py}.
\end{exercise}

\begin{exercise}[$\star$]\label{exo:b2:search-and-games:4}
A tree in which every node has $b \geq 2$ children has $(b^{d+1} - 1)/(b - 1)$ nodes down to depth $d$. Prove it, and
deduce that \omterm{def:b2:search-and-games:iddfs}{iterative deepening} expands at most $b/(b - 1)$ times as many nodes as one depth-limited search to depth $d$,
counting nodes with repetition.
\end{exercise}

\begin{exercise}[$\star\star$]\label{exo:b2:search-and-games:5}
Compute the \omterm{def:b2:search-and-games:minimax}{minimax value} of the tree \texttt{[[[4, 8, 2], [7, 1]], [[3], [9, 6], [5, 0]], [[6, 2], [1, 3, 8]]]} with
\textsc{min} to move at the root, and then with \textsc{max}. Check with \code{exo_minimax.py}.
\end{exercise}

\begin{exercise}[$\star\star$]\label{exo:b2:search-and-games:6}
The function \code{ab_short} of \code{exo_ab_bug.py} starts every node with $\alpha = -\infty$ and $\beta = +\infty$
instead of receiving them from its parent. Is its value right? How many leaves does it evaluate, and why?
\end{exercise}

\begin{exercise}[$\star\star$]\label{exo:b2:search-and-games:7}
In the subtraction game, a move takes 1, 3 or 4 matches, and whoever takes the last wins. Compute the \omterm{def:b2:search-and-games:attractor}{attractor} for the
first player from up to 30 matches, guess the losing \omterm{def:b2:search-and-games:game}{positions}, and prove your guess.
\end{exercise}

\begin{exercise}[$\star\star$]\label{exo:b2:search-and-games:8}
Prove \cref{prop:b2:search-and-games:consistent-admissible}.
\end{exercise}

\begin{exercise}[$\star\star$]\label{exo:b2:search-and-games:9}
Prove that the Manhattan distance of the 8-puzzle is \omterm{def:b2:search-and-games:heuristic}{consistent}, and that it is never smaller than the number of
misplaced tiles. Then compute the true distance of all 181,440 states by one breadth-first search backwards from the
goal, and check both \omterm{def:b2:search-and-games:heuristic}{heuristics} against it.
\end{exercise}

\begin{exercise}[$\star\star\star$]\label{exo:b2:search-and-games:10}
Prove \cref{thm:b2:search-and-games:alpha-beta}.
\end{exercise}

\begin{exercise}[$\star\star\star$]\label{exo:b2:search-and-games:11}
Count the \omterm{def:b2:search-and-games:game}{positions} of noughts and crosses reachable from the empty board, then count them up to the rotations and
reflections of the square, keeping one representative of each class.
\end{exercise}

\begin{exercise}[$\star\star\star$]\label{exo:b2:search-and-games:12}
Prove that \cref{lst:b2:search-and-games:attractor} runs in time $O(|V| + |E|)$, and that a position of the opponent
enters the \omterm{def:b2:search-and-games:attractor}{attractor} exactly when its counter reaches $0$.
\end{exercise}

\section{Problem: The delivery robot and its rival}\label{sec:b2:search-and-games:problem}

\begin{problem}[{Weekend problem --- The delivery robot and its rival}]\label{pb:b2:search-and-games:1}
A delivery robot leaves the Bus Depot with parcels for the Library (place 3), the Museum (12) and the Hospital (8), in any
order. Then a rival robot \omterm{def:b2:text-algorithms:trie}{tries} to catch it on its way to the Library. Last, the robots play noughts and crosses.

\omcode{code/bachelor-2/28-search-and-games/listings/robot_problem.py}{8}{28}{The deliveries as a \omterm{def:b2:search-and-games:state-space}{search problem}: a state
is a place and the deliveries left}\label{lst:b2:search-and-games:problem}

\omoutput{code/bachelor-2/28-search-and-games/outputs/robot_problem.txt}{The best delivery round, the race game solved,
and alpha--beta on noughts and crosses.}\label{out:b2:search-and-games:problem}

\medskip\noindent\textbf{Part I --- The state space.}\par\nopagebreak
\begin{enumerate}
\item How many states does the delivery problem have? Which are goals?
\item Why must the state contain the deliveries left, and not only the place?
\item Why does the route of \cref{out:b2:search-and-games:problem} go through the South Bridge twice? Is that allowed in
the \omterm{def:b2:search-and-games:state-space}{state space}?
\item How many actions does a state have, and what do they cost?
\item Would breadth-first search find the same route? Why not?
\item With $k$ stops instead of 3, how does the number of states grow?
\end{enumerate}

\medskip\noindent\textbf{Part II --- A* on the town.}\par\nopagebreak
\begin{enumerate}[resume, beginpenalty=10000]
\item Prove that $h$, the straight-line distance to the farthest stop left, is \omterm{def:b2:search-and-games:heuristic}{admissible}.
\item Prove that it is \omterm{def:b2:search-and-games:heuristic}{consistent}.
\item Would the sum of the straight-line distances to the stops left be \omterm{def:b2:search-and-games:heuristic}{admissible}? Give a counterexample or a proof.
\item A* expands 56 states and \omterm{def:b2:search-and-games:ucs}{uniform-cost search} 114. Explain where the saving comes from.
\item Propose a better \omterm{def:b2:search-and-games:heuristic}{admissible heuristic}, and say why it is \omterm{def:b2:search-and-games:heuristic}{admissible}.
\item How would A* change if the robot had to return to the depot at the end?
\end{enumerate}

\medskip\noindent\textbf{Part III --- The race, solved backwards.}\par\nopagebreak
\begin{enumerate}[resume, beginpenalty=10000]
\item \omterm{def:b2:search-and-games:game}{Positions} are (robot's place, rival's place, player to move). Check that there are 1,800.
\item Which \omterm{def:b2:search-and-games:game}{positions} form the robot's target, and which the rival's? Why do they get a loop?
\item From the Bus Depot, with the rival at the Station, the robot wins in 3 moves. Describe the play.
\item From the Bus Depot with the rival at the Town Square, neither wins. What do the robots do for ever?
\item From the Allotments, with the rival at the Town Hall, the rival wins. Why can it never be a draw there?
\item The robot wins from 784 \omterm{def:b2:search-and-games:game}{positions} and the rival from 202. Why is the game not symmetric?
\end{enumerate}

\medskip\noindent\textbf{Part IV --- Alpha--beta measured.}\par\nopagebreak
\begin{enumerate}[resume, beginpenalty=10000]
\item On the random tree with $b = 4$ and $d = 8$, compare the leaves evaluated in random order and best first with
$4^8$ and with the formula of \cref{prop:b2:search-and-games:best-case}.
\item In the same time, how much deeper does alpha--beta with perfect ordering search than minimax?
\item On noughts and crosses, alpha--beta finishes 7,330 of the 255,168 games. Why is the saving larger than on the random
trees?
\item Why does alpha--beta call the game with $\alpha = -2$ and $\beta = 2$ rather than infinite bounds?
\item How should a program order the moves of noughts and crosses to approach the best case?
\item State the problem's result: alpha--beta with good ordering visits $L$ of $N$ leaves, near the $b^{d/2}$ bound, with
the values of $L$ and $N$ measured here.
\end{enumerate}
\end{problem}
